\section{Introduction}
\label{sec:intro}

Safety controls that work for one agent need not compose when agents relay
content. An untrusted record can be marked at ingestion, rewritten by an
otherwise honest sender, delivered as a fluent task to a receiver, and acted
on through authority held only by that receiver. The security boundary is
therefore a property of the \emph{agent chain}: what survives transformation,
how it propagates, and which capability can realize the resulting action.
Protocols such as the Model Context Protocol (MCP) and Agent2Agent (A2A) make
such chains easier to build~\cite{mcp-spec-2025,a2a-spec-2025}, but transport
interoperability does not by itself preserve security meaning across derived
messages.

Indirect prompt injection shows that data can steer tool-using
models~\cite{greshake2023not,liu2024formalizing}. Marking, control/data
separation, and capability monitors address different parts of this problem
~\cite{hines2024spotlighting,debenedetti2025camel}. A generative relay sits
between them: it can discard bytes and alter pragmatic emphasis while
preserving the requested action. Evaluating only the first receiver can thus
mistake attenuation for containment; evaluating only executed tool calls can
mistake enforcement for model refusal.

We separate the chain into three observable stages. \emph{Transformation}
records how provenance and action content change during relay.
\emph{Propagation} records whether derived content induces an onward agent
action. \emph{Realization} records whether an attempted action changes
protected state. This decomposition yields three results.

First, literal provenance loss is not equivalent to behavioral invariance.
Across 48 paired
directives, a visible untrusted marker reduced protected-write attempts from
40/48 to 1/48. A normal relay removed the marker from every regenerated text,
yet marked-derived texts produced 22/48 attempts versus 44/48 for
unmarked-derived texts. The paired differences replicated on a byte-identical
Qwen evaluation. Two isolated model raters found that marked-derived texts
contained less actionable content, lower instruction centrality, and stronger
residual provenance cues. The contrast is primarily sender-mediated, and the
post-relay ratings do not identify an independent effect of the missing marker.
An adaptive relay instruction raised broad attempts
from 11/48 to 43/48, showing that the residual is policy-dependent; its exact
protected requests remained descriptively below the unmarked relay (32/48
versus 37/48).

Second, a prospectively selected high-retention stratum supported conditional
two-receiver delegation. Receiver 1 delegated 40/121 messages; receiver 2
delegated 17/40 of those, yielding 17/121 chains with two consecutive
delegation calls. This is a selected-stratum rate, not a population estimate or
a test of selector discrimination. In a separate
fixture, hiding a denied direct-write tool redirected 45/48 elevated requests
to an advertised, non-mutating escalation route. An exact resource policy
prevented unauthorized mutations of a protected canary while models continued
to request writes.

Third, we measure the receiver's response to instruction elevation and weight
it under declared models of elevation and write-authority failure. A covariance
identity isolates coupled failures. These are sensitivity calculations, not
measurements of deployment fault frequency.

Our contributions are paired measurements of how provenance marking changes
relay outputs and receiver actions; an outcome-blind measurement of conditional
two-receiver delegation and action-route migration; and version-pinned
evidence that A2A message derivation can drop supplied lineage. A separate
risk reconstruction relates measured elevation sensitivity to declared
authorization failure topology.
